\documentclass[10pt]{article}
%-*-*-*-*-*-*-*-*-*-*-*-*-*-*-*-*-*-*-*Packages--*--*--*--*--*--*--*--*--*--*--*--*--*
\usepackage{amsmath}
\usepackage{tikz}
\usepackage{epigraph}
\usepackage{lipsum}
\usepackage{hyperref} 
\usepackage{listings}
\usepackage{fancyhdr}
 \usepackage[utf8]{inputenc}
\usepackage{hyperref} 
\usepackage[spanish]{babel}
%\usepackage{simpsons}
%\usepackage[dvips]{graphicx}
\usepackage{amsmath}
\usepackage{amsthm}
%\usepackage{textcomp}
\usepackage{amssymb}
\usepackage{latexsym}
\usepackage{graphicx}
%\usepackage[ansinew]{inputenc}
\usepackage{color}
%\usepackage{pstricks, pst-node, pst-plot, pst-circ}
%\usepackage{moredefs}
%\usepackage{pstricks}
%\usepackage{pst-circ}
%\usepackege{pst-node}
%\usepackage{pst-plot}
\usepackage{moredefs}
%\usepackage{mcode}
%\usepackage[above]{placeins}
\usepackage{fancybox}
\usepackage{subfig}
\usepackage{float}
%\usepackage{mcode} para colorear codigo de matlab
\usepackage{xcolor}
\usepackage{wallpaper}
\usepackage{textcomp}
%-*-*-*-*-*-*-*-*-*-*-*-*-*-*-*-*-*-*--*--*--*--*--*--*--*--*--*--*--*--*--*--*

%*****************************definiciones*******************************
\oddsidemargin -0.2in \textwidth 7.0in \topmargin -0.9in \textheight
9.0in
\parindent 3em
\hyphenation{e-jem-plo
e-rro-res de-pen-dien-te co-rrien-te res-pues-ta fi-gu-ran
cons-tan-tes o-pe-ra-cion i-lu-mi-na-cion}

\pagestyle{fancy}
\headheight 50pt 
\renewcommand{\arrayrulewidth}{3pt} 
\newtheorem{defi}{{\it Definición}}[section]
%@@@@@@@@@@@@@@@@@@@@@@@@@@@cabecera y pie de pagina @@@@@@@@@@@@@@@@@@@@@@@@@@@@@@@@@@@@@@@@@@@@@@@@@@@@@@

\rhead{\includegraphics[width=.075\textwidth]{logo_iaci2.eps}}
\chead{}
\lhead{Teoría de las Telecomunicaciones}
\rfoot{1º cuatrimestre 2015  }
\cfoot{\thepage}
\lfoot{\bf{Universidad Nacional de Quilmes}}
%@@@@@@@@@@@@@@@@@@@@@@@@@@@@@@@@@@@@@@@@@@@@@@@@@@@@@@@@@@@@@@@@@@@@@@@@@@@@@@@@@@@@@@@@@@@@@@@@@@@@@@@@@@@

%::::::::::::::::::::::Colores varios:::::::::::::::::::::::::::::::::::::::::::::::::::::::::::::::::::::::

\definecolor{azul-claro}{rgb}{0.335,0.89,1}
\definecolor{gainsboro}{rgb}{0.86,0.86,0.86}
\definecolor{mediumseagreen}{rgb}{0.24,0.7,0.44}
\definecolor{moccasin}{rgb}{1,0.89,0.71}
\definecolor{cornflowerblue}{rgb}{0.39,0.58,0.93}
\definecolor{lightgray}{rgb}{0.83,0.83,0.83}
\definecolor{darkgrey}{rgb}{0.66,0.66,0.66}
\definecolor{darkslategray}{rgb}{0.18,0.31,0.31}
\definecolor{lavender}{rgb}{0.9,0.9,0.98}
\definecolor{azure}{rgb}{0.94,1,1}
\definecolor{honeydew}{rgb}{0.94,1,0.94}

%:::::::::::::::::::::::::::::::::::::::::::::::::::::::::::::::::::::::::::::::::::::::::::::::::::::::::::::
%Requiere \usepackage{xcolor}
\renewcommand{\arrayrulewidth}{1pt} 
\newcommand{\mcaja}[1]{%
{{\fboxsep 10pt \fboxrule 1pt%
\fcolorbox{black}{orange}{%
\color{black} \Huge #1}}}
}
\newcommand{\nuevobox}[1]{%
{{\fboxsep 14pt \fboxrule 1pt%
\fcolorbox{black}{darkgrey}{%
\color{lavender} \huge #1}}}
}
\newcommand{\ssection}[1]{\section[#1]{\mcaja{#1}}}
\makeatletter
\newcommand{\sect}[1]{\subsection[#1]{\nuevobox{#1}}}
\makeatletter
\def\section{\@ifstar\unnumberedsection\numberedsection}
\def\numberedsection{\@ifnextchar[%]
\numberedsectionwithtwoarguments\numberedsectionwithoneargument}
\def\unnumberedsection{\@ifnextchar[%]
\unnumberedsectionwithtwoarguments\unnumberedsectionwithoneargument}
\def\numberedsectionwithoneargument#1{\numberedsectionwithtwoarguments[#1]{#1}}
\def\unnumberedsectionwithoneargument#1{\unnumberedsectionwithtwoarguments[#1]{#1}}
\def\numberedsectionwithtwoarguments[#1]#2{%
\ifhmode\par\fi
\removelastskip
\vskip 3ex\goodbreak
\refstepcounter{section}%
\begingroup
%\noindent
\leavevmode\large\bfseries\raggedright\mcaja%%
\thesection\ #2\par\nobreak
\endgroup
\noindent\hrulefill\nobreak
\vskip 2ex\nobreak
\addcontentsline{toc}{section}{%
\protect\numberline{\thesection}%
#1}%
}
\def\unnumberedsectionwithtwoarguments[#1]#2{%
\ifhmode\par\fi
\removelastskip
\vskip 3ex\goodbreak
% \refstepcounter{section}%
\begingroup
\noindent
\leavevmode\Large\bfseries\raggedright
% \thesection\
#2\par\nobreak
\endgroup
\noindent\hrulefill\nobreak
\vskip 0ex\nobreak
\addcontentsline{toc}{section}{%
%
\protect\numberline{\thesection}%
#1}%
}
\makeatother
%%%Cap\’itulos
\usepackage{helvet}
%\usepackage{psboxit,pstcol}
\makeatletter
\def\@makechapterhead#1{%
{\parindent \z@ \raggedright \reset@font
\hbox to \hsize{%
\rlap{\raisebox{-2.5em}{\raisebox{\depth}{%%% Necesita la imagen "imgCapitulo"
\includegraphics[width=10em]{logo_iaci2.eps}}}}%
\rlap{\hbox to 6em{\hss
\reset@font\sffamily\fontsize{8em}{8em}\selectfont\black
\thechapter\hss}}%
\hspace{10em}%
\vbox{%
\advance\hsize by -10em
\reset@font\fontfamily{hv}\bfseries\Large
#1
\par
}%
}}%
\vskip 5pt
\hrulefill
\vskip 50pt
}
\makeatother
%----------------------------------------------------------------------------------------
% Theorems
\newtheorem{theorem}{Teorema}[section]
\newtheorem{corollary}{Corolario}[theorem]
\newtheorem{lemma}[theorem]{Lema}
%----------------------------------------------------------------------------------------


%----------------------------------------------------------------------------------------
%	TITLE PAGE
%----------------------------------------------------------------------------------------
\definecolor{titlepagecolor}{cmyk}{1,.73,0,.37}


\newcommand\titlepagedecoration{%
\begin{tikzpicture}[remember picture,overlay,shorten >= -10pt]

\coordinate (aux1) at ([yshift=-15pt]current page.north east);
\coordinate (aux2) at ([yshift=-410pt]current page.north east);
\coordinate (aux3) at ([xshift=-4.5cm]current page.north east);
\coordinate (aux4) at ([yshift=-150pt]current page.north east);

\begin{scope}[titlepagecolor!40,line width=12pt,rounded corners=12pt]
\draw
  (aux1) -- coordinate (a)
  ++(225:5) --
  ++(-45:5.1) coordinate (b);
\draw[shorten <= -10pt]
  (aux3) --
  (a) --
  (aux1);
\draw[opacity=0.6,titlepagecolor,shorten <= -10pt]
  (b) --
  ++(225:2.2) --
  ++(-45:2.2);
\end{scope}
\draw[titlepagecolor,line width=8pt,rounded corners=8pt,shorten <= -10pt]
  (aux4) --
  ++(225:0.8) --
  ++(-45:0.8);
\begin{scope}[titlepagecolor!70,line width=6pt,rounded corners=8pt]
\draw[shorten <= -10pt]
  (aux2) --
  ++(225:3) coordinate[pos=0.45] (c) --
  ++(-45:3.1);
\draw
  (aux2) --
  (c) --
  ++(135:2.5) --
  ++(45:2.5) --
  ++(-45:2.5) coordinate[pos=0.3] (d);   
\draw 
  (d) -- +(45:1);
\end{scope}
\end{tikzpicture}%
}


\newcommand*{\plogo}{\fbox{$\mathcal{PL}$}}
\newcommand*{\titleGP}{\begingroup % Create the command for including the title page in the document
\centering % Center all text
\vspace*{\baselineskip} % White space at the top of the page

\rule{\textwidth}{1.6pt}\vspace*{-\baselineskip}\vspace*{2pt} % Thick horizontal line
\rule{\textwidth}{0.4pt}\\[\baselineskip] % Thin horizontal line

{\LARGE Teoría de las telecomunicaciones\\[0.3\baselineskip] Laboratorio 1}\\[0.2\baselineskip] % Title

\rule{\textwidth}{0.4pt}\vspace*{-\baselineskip}\vspace{3.2pt} % Thin horizontal line
\rule{\textwidth}{1.6pt}\\[\baselineskip] % Thick horizontal line

\scshape % Small caps
 Martín Noblía 
\vspace*{2\baselineskip} % Whitespace between location/year and editors

Profesores: \\[\baselineskip]
{\Large Fabian Iakinchuk\\ Martín Castillo\par}
\vfill % Whitespace between editor names and publisher logo_iaci2
\includegraphics[width=.175\textwidth]{logo_iaci2.eps}

% Editor list
{\itshape Universidad Nacional de Quilmes \par} % Editor affiliation

\vfill % Whitespace between editor names and publisher logo

%\plogo \\[0.3\baselineskip] % Publisher logo
{\scshape 2015} \\[0.3\baselineskip] % Year published
%{\large THE PUBLISHER}\par % Publisher

\endgroup}

%----------------------------------------------------------------------------------------
%	BLANK DOCUMENT
%----------------------------------------------------------------------------------------


\begin{document}
\pagestyle{empty} 

\titleGP
\titlepagedecoration

\newpage
\pagestyle{fancy}

\section{Problema 1}

Dada una fuente de datos en un archivo de texto nombrado:\textit{texto.txt}.
Escribir un algoritmo de manera tal que calcule la probabilidad de ocurrencia 
de cada símbolo existente en la fuente.
Mostrar en un listado (sea en pantalla o en otro archivo \textit{.txt} cada 
elemento de la fuente con su probabilidad calculada asociada.
Calcular la entropía de la fuente.
\subsection{resolución}
Para resolver este problema utilizamos el lenguaje de programación Julia(\url{http://www.julialang.org})
el cual es \textit{opensource}, dinámico y tiene velocidades de procesamiento relativas a C entre otras 
caracteristicas.

El script que realiza la tarea es el siguiente:

\begin{verbatim}
#= Script para el problema 1
desde consola:
julia problema1.jl
=#
# incluimos el modulo
include("teleco.jl")
# leemos el archivo
T = open(readall,"../Archivos/texto.txt");
# calculamos las frecuencias y probabilidades 
d,p = Teleco.probs(T[1:end-1]);
# Mostramos los resultados
Teleco.show_source(d)
# Calculamos la entropia de la fuente
h = Teleco.entropy(p);
print("La fuente discreta texto.txt tiene una entropia: $h [bits/simbolo]")
\end{verbatim}

Las definiciones de funciones estan agrupadas en el modulo \textit{Teleco}
que esta definido en el archivo \textit{teleco.jl}. El cual mostramos a continuación:

\begin{verbatim}
# Modulo para el laboratorio1:Teoria de las telecomunicaciones
module Teleco
    

function probs(iter)
    # Funcion para calcular las probabilidades de ocurrencia de cada simbolo
    # en una fuente de informacion discreta sin memoria.
    # Input:
    # -----
    # iter: Cualquier fuente iterable(texto,vectores...etc) 
    # Output:
    # ------
    # d: Diccionario con los elementos de la fuente y sus frecuencias
    # p: Vector de probabilidades

    d = Dict{eltype(iter), Int}()
    l = length(iter)
    
    for e in iter
        d[e] = get(d, e, 0) + 1 
    end
    p = collect((values(d))) / l 
    return d,p
end

function freqs_pr(iter)
    # Funcion para calcular 
    pr = Collections.PriorityQueue()
    
    for e in iter
        pr["$e"] = get(pr, "$e", 0) + 1 
    end
    
    return pr
end


function entropy(p)
   # Funcion para calcular la entropia de una 
   # fuente discreta sin memoria
   # Input:
   # -----
   # p: Vector de probabilidades de la fuente
   # Output:
   # ------
   # h: Entropia de la fuente
    i = log2(1./p)
    aux = p.*i
    h = sum(aux)
    
    return h

end

function show_source(d)
    # Funcion para mostrar los elementos de una fuente 
    # y su frecuencia de aparicion
    for (v,k) in zip(values(d),keys(d))
        println("elemento: $k --> frecuencia: $v")
    end
end    

export probs, probs_pr,entropy, show_source


end

\end{verbatim}

Cuya salida es:

\begin{verbatim}
elsuizo@debian:~/Telecomunicaciones/Laboratorios/Lab1/Programas/Julia$ ~/julia/./julia problema1.jl 
elemento:   --> frecuencia: 22
elemento: á --> frecuencia: 4
elemento: l --> frecuencia: 6
elemento: m --> frecuencia: 2
elemento: r --> frecuencia: 5
elemento: f --> frecuencia: 1
elemento: o --> frecuencia: 3
elemento: t --> frecuencia: 2
elemento: a --> frecuencia: 14
elemento: E --> frecuencia: 1
elemento: v --> frecuencia: 2
elemento: n --> frecuencia: 5
elemento: e --> frecuencia: 13
elemento: w --> frecuencia: 2
elemento: h --> frecuencia: 3
elemento: d --> frecuencia: 3
elemento: c --> frecuencia: 2
elemento: ; --> frecuencia: 1
elemento: p --> frecuencia: 2
elemento: i --> frecuencia: 4
elemento: q --> frecuencia: 5
elemento: s --> frecuencia: 4
elemento: u --> frecuencia: 8
La fuente discreta texto.txt tiene una entropia: 3.989736156368347 [bits/simbolo]
\end{verbatim}

\section{Problema 2}

Dado un archivo de texto nombrado “texto2.txt” el cual posee una longitud fija de
30 caracteres ASCII:

Escribir un algoritmo que codifique
cada elemento de la fuente usando codificación Huffman. (Dicha codificación puede ser
mostrada en pantalla o en un archivo de salida *.txt).
El software creado debe calcular la entropía y la longitud media del código generado.
Calcular en forma manual o con el mismo software la eficiencia de la compresión.
Verifique si el código generado es óptimo.

\subsection{resolución}


\newpage

\section{Problema 3}
Dado dos archivos de texto llamados Castellano.txt e ingles.txt, los cuales están en
diferente lenguaje, comprimir los mismos en formato ZIP y:
\begin{itemize}
\item a) Realizar una tabla donde se pueda visualizar el cociente entre: (tamaño nuevo /
tamaño original) x100.
\item b) Realizar una tabla con la probabilidad de ocurrencia de cada uno de los caracteres
para cada texto.
\item c) Analizar los resultados de a) y b). Sacar conclusiones.
\end{itemize}


\newpage

\begin{thebibliography}{6}
 \bibitem{Principal}{Preparación y evaluación de Proyectos,cuarta edición,Nassir Sapag Chain, Reinaldo Sapag Chain}\\
 \textbf{Enlaces}:
\begin{center}
\url{www.google.com.ar/}\\
\url{http://www.lanacion.com.ar/}\\
\url{http://www.espacioyconfort.com.ar/ }\\
\end{center}

\end{thebibliography}



\end{document}
